\documentclass[14pt,aspectratio=1610]{beamer}
%\documentclass[14pt,aspectratio=1610,handout]{beamer}
\usepackage{csu-theme}
\usepackage{forest}
\usepackage{booktabs}

\title{\Huge{} Quick Sort \&\\Merge Sort}
\subtitle{CSCI 315: Data Structure Analysis\\Chapter 18}
%\author{Sean T. Hayes, Ph.D.}
\institute{Charleston Southern University}

%% Comment out date to use default (today)
% \date{February 18, 2025}
\subject{Software Programming}
%\keywords{}

\setlength{\parskip}{0.5em}


\definecolor{MyGreen}{HTML}{61ff7e}
\usepackage{linked-lists}
\usepackage{arrays}


\tikzset{
  node of list/.append style = {
    minimum height=6mm,
    minimum width=3mm,
    node distance=6mm,
    font=\small,
  },
  link pointer of list/.append style= {
    minimum height=6mm,
    minimum width=3mm
  },
  pointer/.append style = {
    minimum height=4mm,
    minimum width=4mm,
    node distance=1pt
  },
  label of pointer/.append style = {
    font=\ttfamily\small
  },
  node of array/.append style = {
     minimum height=8mm,
     minimum width=10.8mm,
     font=\large
  },
  label of array el/.append style={
    % rotate=-67.5,
    % anchor=west,
    % align=right,
    % yshift=-3mm,
    font=\sffamily,
    % rotate=-12,
  },
}

\begin{document}

\begin{frame}
  \titlepage
\end{frame}

\section{Best Performance}

\begin{frame}{Lower Bound on Comparison-Based Sorting}
\center\Large
Can we do better than $O(n^2)$ for sorting algorithms?\vspace{1em}

What is the best we can do?
\end{frame}

\begin{frame}{Comparison Tree}
\large
\begin{itemize}
\item
  \emph{Comparison tree}: graph used to trace the execution of a comparison-based algorithm.
  \begin{itemize}
  \item
    Let $L$ be a list of $n$ distinct elements, where $n > 0$\\
    For any $j$ and $k$, where $1 \le j \le n$,  $1 \le k \le n$,\\
    either $L[j] \le L[k]$  or  $L[j] > L[k] $.
  \end{itemize}
\item
  \emph{Binary tree}: each comparison has two outcomes.
\end{itemize}
\end{frame}

\begin{frame}{Terminology of Comparison Trees}
\begin{columns}
  \column{0.66\textwidth}
\begin{itemize}

\item
  \emph{Node}: represents a comparison
  \begin{itemize}
  \item
    Labeled as $j:k$ (comparison of $L[j]$ with $L[k]$)
  \item
    If $L[j] < L[k]$, follow the left branch; otherwise, follow the right branch.
  \end{itemize}
\item
  \emph{Leaf}: represents the final ordering of the nodes
\item
  \emph{Root}: the top node
\item
  \emph{Branch}: line that connects two nodes
\item
  \emph{Path}: sequence of branches from one node to another
\end{itemize}
\column{0.33\textwidth}
\begin{center}
  \begin{forest}
  for tree = {
    circle, draw,
    inner sep=3pt,
    color=CSUblue,
    fill=white,
    fit=tight,
    edge=->,
    edge=very thick
  }
    [A
      [B
        [D
          [G]
          [, phantom]
        ]
        [E]
      ]
      [C
        [F]
        [, phantom]
      ]
    ]
  \end{forest}
\end{center}
\end{columns}
\end{frame}

\begin{frame}{Comparison Tree for Sorting 3 Items}
\center{}
\begin{forest}
for tree = {
  circle, draw,
  color=CSUblue,
  fill=white,
  %fit=band,
  l sep=2.6cm,
  s sep+=3mm,
  edge=->,
  edge=very thick,
  grow=east,
  for leaves={rectangle, fill=CSUgold},
  font=\small
}
[{$L_1:L_2$} ,name=root,
    [{$L_2:L_3$}, edge label = {node [midway,below,sloped] {$L_1 \le L_2$} }
      [{$\mathbf{1, 2, 3}$}, edge label = {node [midway,below,sloped] {$L_2 \le L_3$} }]
      [{$L_1:L_3$}, edge label = {node [midway,above,sloped] {$L_2 > L_3$} }
        [{$\mathbf{1, 3, 2}$}, edge label = {node [midway,below,sloped] {$L_1 \le L_3$} }]
        [{$\mathbf{3, 1, 2}$}, edge label = {node [midway,above,sloped] {$L_1 > L_3$} }]
      ]
    ]
    [{$L_2:L_3$}, edge label = {node [midway,above,sloped] {$L_1 > L_2$} }
      [{$L_1:L_3$}, edge label = {node [midway,below,sloped] {$L_2 \le L_3$} }
        [{$\mathbf{2, 1, 3}$}, edge label = {node [midway,below,sloped] {$L_1 \le L_3$} }]
        [{$\mathbf{2, 3, 1}$}, edge label = {node [midway,above,sloped] {$L_1 > L_3$} }]
      ]
      [{$\mathbf{3, 2, 1}$}, edge label = {node [midway,above,sloped] {$L_2 > L_3$} }]
    ]
]
\end{forest}
\end{frame}

\begin{frame}{Comparison Tree}
\begin{itemize}
\item<+->
  A unique permutation of the elements of $L$ is associated with each root-to-leaf path.
  \begin{itemize}
  \item
    Because the sort algorithm only moves the data and makes comparisons.
  \end{itemize}
\item<+->
  For a list of $n$ elements, $n > 0$, there are $n!$ different permutations
  \begin{itemize}
  \item
    Any of these might be the correct ordering of $L$.
  \end{itemize}
\item<+->
  Thus, the tree must have \emph{at least} $n!$ leaves.
\end{itemize}
\end{frame}

\begin{frame}{Lower Bound on Comparison-Based Sorting}
\large
Theorem: Any sorting algorithm that sorts a list of $n$ elements (by comparison of only the keys) makes \emph{at least} $O(n\log_2n)$ key comparisons in its worst case.\vspace{1em}

\pause{}
We won't prove this theorem, but unlike the previous algorithms discussed, there are algorithms that, on average, are $O(n\log_2n)$.
\end{frame}

\section{Quick Sort}

\begin{frame}{Introduction to Quick Sort}
\emph{Quick sort}: uses the divide-and-conquer technique.

\begin{itemize}
\item
  The list is \textit{partitioned} into two sub-lists.
\item
  Each sub-list is then sorted.
\item
  Sorted sub-lists are combined into one list in such a way that the combined
  list is sorted.
\item
  All the sorting work occurs when partitioning the list.
\end{itemize}
\end{frame}

\begin{frame}{Pivot Element}
Choose an element to split the list \emph{pivot element}.
\begin{itemize}
\item
  Move elements less than the pivot value to the lower sub-list.
\item
  Move elements greater than the pivot value to the upper sub-list.
\item
  The pivot can be chosen in several ways.\\
  Ideally, the pivot will divide the list into two sub-lists of near-equal size.
\end{itemize}
\end{frame}

\section{Quick Sort: Example}

\begin{frame}[t]{Quick Sort Example}
  \vspace{2.1em}
  \centerline{%
  \begin{tikzpicture}[]
    \ArrayLabel{}
    \alt<1>{
      \ArrayNode{52}
    }{
      \ArrayNode[MyGreen]{52}
      \ArrayElLabel[val]{pivot}
    }
    \ArrayList{55, 87, 13, 78, 96, 32, 48, 22, 11, 58, 66, 88, 45}

    \ArrayGroup{el52}{el45}{Unsorted List}
  \end{tikzpicture}}

  \vspace{0.7em}
  \centerline{%
  \begin{tikzpicture}
    \ArrayLabel{}
    \ArrayList[CSUgold]{45, 13, 32, 48, 22, 11}
  
    \ArrayNode[MyGreen]{52}
    \ArrayElLabel[val]{pivot}
  
    \ArrayList[MyPurple]{55, 78, 96, 58, 66, 88, 87}
  
    \ArrayGroup{el45}{el11}{Lower}
    \ArrayGroup{el55}{el87}{Upper}
  \end{tikzpicture}}
\end{frame}

\begin{frame}[b]{Quick Sort Example}
\centerline{%
\begin{tikzpicture}[]
  \ArrayLabel{}
  \ArrayNode[MyGreen]{52}
  \ArrayElLabel<-12mm><0mm>[val]{pivot}
  \onslide<-4>{
      \ArrayElLabel<-8mm><8mm>[val]{small}
  }
  \alt<1>{
    \ArrayNode{55}
  }{
    \alt<-5>{
      \alt<-4>{
        \ArrayNode[MyPurple]{55}
      }{
        \ArrayNode[CSUgold]{55}
      }

    }{
      \ArrayNode[CSUgold]{13}
    }
    \onslide<5->{
      \ArrayElLabel<-8mm><0mm>[val]{small}
    }
  }

  \alt<-2>{
    \ArrayNode{87}
  }{
    \ArrayNode[MyPurple]{87}
  }

  \alt<-5>{
    \ArrayNode{13}
  }{
    \ArrayNode{55}
  }
  \ArrayList{78, 96, 32, 48, 22, 11, 58, 66, 88, 45}

  \onslide<1>{
    \ArrayGroup{el52}{el45}{Unsorted List}
  }
  \alt<1>{
    \ArrayElLabel<-4mm><4mm>[el55]{index}
  }{
    \alt<-2>{
      \ArrayElLabel[el87]{index}
      \ArrayGroup{el55}{el55}{Upper}
    }{

      \alt<-4>{
        \ArrayElLabel<-8mm><0mm>[el13]{index}
        \ArrayGroup{el55}{el87}{Upper}
      }{
        \alt<-5>{
          \ArrayElLabel<-8mm><0mm>[el13]{index}
          \ArrayGroup{el55}{el55}{Lower}
        }{
          \ArrayElLabel<-8mm><0mm>[el55]{index}
          \ArrayGroup{el13}{el13}{Lower}
        }
        \ArrayGroup<16pt>{el87}{el87}{Upper}
      }
    }
  }
  \only<4-5>{
    \ArraySwap{el55}{el13}
  }
  \only<6>{
    \ArraySwap{el13}{el55}
  }
  \onslide<7>{}

  % reserve the space for this in a later slide.
  \invisible{%
    \ArraySwap{el32}{el45}%
  }
\end{tikzpicture}}
\vspace{3.5em}
\end{frame}

\begin{frame}[b]{Quick Sort Example}

\centerline{%
\begin{tikzpicture}
  \ArrayLabel{}
  \ArrayNode[MyGreen]{52}
  \ArrayElLabel<-12mm><0mm>[val]{pivot}
  \ArrayNode[CSUgold]{13}
  \ArrayElLabel<-8mm><0mm>[val]{small}
  \ArrayList[MyPurple]{87, 55}
  \alt<1>{
    \ArrayNode[]{78}
  } {
    \ArrayNode[MyPurple]{78}
  }

  \alt<-2>{
    \ArrayNode{96}
  }{
    \ArrayNode[MyPurple]{96}
  }
  \ArrayList{32, 48, 22, 11, 58, 66, 88, 45}

  \ArrayGroup{el13}{el13}{Lower}
  \alt<1> {
    \ArrayElLabel[el78]{index}

    \ArrayGroup{el87}{el55}{Upper}
  }{
    \alt<-2>{
      \ArrayElLabel[el96]{index}
      \ArrayGroup{el87}{el78}{Upper}
    }{
      \ArrayElLabel[el32]{index}
      \ArrayGroup{el87}{el96}{Upper}
    }
  }
  % reserve the space for this in a later slide.
  \invisible{%
    \ArraySwap{el32}{el45}%
  }
  \onslide<3>{}
\end{tikzpicture}}
\vspace{3.5em}
\end{frame}

\begin{frame}[b]{Quick Sort Example}
\centerline{%
\begin{tikzpicture}
  \ArrayLabel{}
  \ArrayNode[MyGreen]{52}
  \ArrayElLabel<-12mm><0mm>[val]{pivot}
  \ArrayNode[CSUgold]{13}

  \alt<1>{
    \ArrayElLabel<-8mm><0mm>[val]{small}
    \ArrayNode[MyPurple]{87}
  }{
    \ArrayNode[CSUgold]{87}
    \ArrayElLabel{small}
  }

  \ArrayList[MyPurple]{55, 78, 96}
  \ArrayNode[]{32}
  \ArrayList{48, 22, 11, 58, 66, 88, 45}

  \ArrayElLabel[el32]{index}


  \ArraySwap{el87}{el32}

  \alt<1> {
    \ArrayGroup{el13}{el13}{Lower}
    \ArrayGroup{el87}{el96}{Upper}
  }{
    \ArrayGroup{el13}{el87}{Lower}
    \ArrayGroup{el55}{el96}{Upper}
  }

  % reserve the space for this in a later slide.
  \invisible{%
    \ArraySwap{el32}{el45}%
  }

  \onslide<2>{}
\end{tikzpicture}}
\vspace{3.5em}
\end{frame}


\begin{frame}[b]{Quick Sort Example}
\centerline{%
\begin{tikzpicture}
  \ArrayLabel{}
  \ArrayNode[MyGreen]{52}
  \ArrayElLabel<-12mm><0mm>[val]{pivot}
  \ArrayList[CSUgold]{13, 32}
  \ArrayElLabel[val]{small}

  \ArrayList[MyPurple]{55, 78, 96}
  \alt<-2>{
    \ArrayNode{87}
  }{
    \ArrayNode[MyPurple]{87}
  }
  \ArrayList{48, 22, 11, 58, 66, 88, 45}

  \ArrayGroup{el13}{el32}{Lower}
  \alt<-2>{
    \ArrayElLabel[el87]{index}
    \ArrayGroup{el55}{el96}{Upper}
  }{
    \ArrayElLabel[el48]{index}
    \ArrayGroup{el55}{el87}{Upper}
  }

  \onslide<1> {
    \ArraySwap{el32}{el87}
  }

  % reserve the space for this in a later slide.
  \invisible{%
    \ArraySwap{el87}{el45}%
  }

  \onslide<3> {}
\end{tikzpicture}}
\vspace{3.5em}
\end{frame}

\begin{frame}[b]{Quick Sort Example}
\centerline{%
\begin{tikzpicture}
  \ArrayLabel{}
  \ArrayNode[MyGreen]{52}
  \ArrayElLabel<-12mm><0mm>[val]{pivot}
  \ArrayList[CSUgold]{13, 32}

  \alt<1>{
    \ArrayNode[MyPurple]{55}
  }{
    \alt<-2>{
      \ArrayNode[CSUgold]{55}
    }{
      \ArrayNode[CSUgold]{48}
    }
  }
  \ArrayList[MyPurple]{78, 96, 87}
  \alt<-2>{
    \ArrayNode{48}
  }{
    \alt<-4>{
      \ArrayNode{55}
    }{
      \ArrayNode[MyPurple]{55}
    }
  }
  \ArrayList{22, 11, 58, 66, 88, 45}

  \onslide<1-2> {
    \ArraySwap{el55}{el48}
  }
  \onslide<3> {
    \ArraySwap{el48}{el55}
  }

  \alt<-2>{
    \ArrayElLabel[el48]{index}
  }{
    \alt<-4>{
      \ArrayElLabel[el55]{index}
    }{
      \ArrayElLabel[el22]{index}
    }
  }

  \alt<1> {
    \ArrayElLabel[el32]{small}
    \ArrayGroup{el13}{el32}{Lower}
    \ArrayGroup{el55}{el87}{Upper}
  }{
    \alt<-2>{
      \ArrayElLabel[el55]{small}
      \ArrayGroup{el13}{el55}{Lower}
    }{
      \ArrayElLabel[el48]{small}
      \ArrayGroup{el13}{el48}{Lower}
    }
    \alt<-4>{
      \ArrayGroup{el78}{el87}{Upper}
    }{
      \ArrayGroup{el78}{el55}{Upper}
    }
  }

  % reserve the space for this in a later slide.
  \invisible{%
    \ArraySwap{el87}{el45}%
  }

  \onslide<5>{}
\end{tikzpicture}}
\vspace{3.5em}
\end{frame}

\begin{frame}[b]{Quick Sort Example}
\centerline{%
\begin{tikzpicture}
  \ArrayLabel{}
  \ArrayNode[MyGreen]{52}
  \ArrayElLabel<-12mm><0mm>[val]{pivot}
  \ArrayList[CSUgold]{13, 32, 48}

  \alt<1>{
    \ArrayElLabel[val]{small}
    \ArrayNode[MyPurple]{78}
  }{
    \alt<-2>{
      \ArrayNode[CSUgold]{78}
    }{
      \ArrayNode[CSUgold]{22}
    }
    \ArrayElLabel[val]{small}
  }

  \ArrayList[MyPurple]{96, 87, 55}

  \alt<-2>{
    \ArrayNode{22}
  }{
    \alt<-4>{
      \ArrayNode{78}
    } {
      \ArrayNode[MyPurple]{78}
    }
  }

  \alt<-4>{
    \ArrayElLabel{index}
    \ArrayNode{11}
  }{
    \ArrayNode{11}
    \ArrayElLabel{index}
  }

  \ArrayList{58, 66, 88, 45}

  \onslide<1-2> {
    \ArraySwap{el78}{el22}
  }
  \onslide<3> {
    \ArraySwap{el22}{el78}
  }

  \alt<1> {
    \ArrayGroup{el13}{el48}{Lower}
    \ArrayGroup{el78}{el55}{Upper}
  }{
    \alt<-2>{
      \ArrayGroup{el13}{el78}{Lower}
    }{
      \ArrayGroup{el13}{el22}{Lower}
    }
    \alt<-4>{
      \ArrayGroup{el96}{el55}{Upper}
    }{
      \ArrayGroup{el96}{el78}{Upper}
    }
  }

  % reserve the space for this in a later slide.
  \invisible{%
    \ArraySwap{el87}{el45}%
  }

  \onslide<5>{}
\end{tikzpicture}}
\vspace{3.5em}
\end{frame}

\begin{frame}[b]{Quick Sort Example}
\centerline{%
\begin{tikzpicture}
  \ArrayLabel{}
  \ArrayNode[MyGreen]{52}
  \ArrayElLabel<-12mm><0mm>[val]{pivot}
  \ArrayList[CSUgold]{13, 32, 48, 22}

  \alt<1>{
    \ArrayElLabel[val]{small}
    \ArrayNode[MyPurple]{96}
  }{
    \alt<-2>{
      \ArrayNode[CSUgold]{96}
    }{
      \ArrayNode[CSUgold]{11}
    }
    \ArrayElLabel[val]{small}
  }

  \ArrayList[MyPurple]{87, 55, 78}

  \alt<-2>{
    \ArrayNode{11}
  }{
    \alt<-4>{
      \ArrayNode{96}
    } {
      \ArrayNode[MyPurple]{96}
    }
  }

  \alt<-4>{
    \ArrayElLabel{index}
    \ArrayNode{58}
  }{
    \ArrayNode{58}
    \ArrayElLabel{index}
  }

  \ArrayList{66, 88, 45}

  \onslide<1-2> {
    \ArraySwap{el96}{el11}
  }
  \onslide<3> {
    \ArraySwap{el11}{el96}
  }

  \alt<1> {
    \ArrayGroup{el13}{el22}{Lower}
    \ArrayGroup{el96}{el78}{Upper}
  }{
    \alt<-2>{
      \ArrayGroup{el13}{el96}{Lower}
    }{
      \ArrayGroup{el13}{el11}{Lower}
    }
    \alt<-4>{
      \ArrayGroup{el87}{el78}{Upper}
    }{
      \ArrayGroup{el87}{el96}{Upper}
    }
  }

  % reserve the space for this in a later slide.
  \invisible{%
    \ArraySwap{el87}{el45}%
  }

  \onslide<5>{}
\end{tikzpicture}}
\vspace{3.5em}
\end{frame}


\begin{frame}[b]{Quick Sort Example}
\centerline{%
\begin{tikzpicture}
  \ArrayLabel{}
  \ArrayNode[MyGreen]{52}
  \ArrayElLabel<-12mm><0mm>[val]{pivot}
  \ArrayList[CSUgold]{13, 32, 48, 22, 11}
  \ArrayElLabel[val]{small}

  \ArrayList[MyPurple]{87, 55, 78, 96, 58}
  \alt<1> {
    \ArrayNode[]{66}
  } {
    % \ArrayNode[MyPurple]{66}
  }
  \ArrayList{88, 45}

  \ArrayElLabel[el66]{index}

  \onslide<1> {
    \ArrayGroup{el13}{el11}{Lower}
    \ArrayGroup{el87}{el58}{Upper}
  }
  % \onslide<2> {
  %   \ArrayGroup{el13}{el11}{Lower}
  %   \ArrayGroup{el87}{el66}{Upper}
  % }

  % reserve the space for this in a later slide.
  \invisible{%
    \ArraySwap{el87}{el45}%
  }

\end{tikzpicture}}
\vspace{3.5em}
\end{frame}

\begin{frame}[b]{Quick Sort Example}
\centerline{%
\begin{tikzpicture}
  \ArrayLabel{}
  \ArrayNode[MyGreen]{52}
  \ArrayElLabel<-12mm><0mm>[val]{pivot}
  \ArrayList[CSUgold]{13, 32, 48, 22, 11}
  \ArrayElLabel{small}

  \ArrayList[MyPurple]{87, 55, 78, 96, 58, 66}
  \alt<1> {
    \ArrayNode[]{88}
    \ArrayElLabel[el88]{index}
  } {
    \ArrayNode[MyPurple]{88}
  }
  \ArrayList{45}

  \onslide<2->{
    \ArrayElLabel[el45]{index}
  }

  \ArrayGroup{el13}{el11}{Lower}
  \alt<1> {
    \ArrayGroup{el87}{el66}{Upper}
  }{
    \ArrayGroup{el87}{el88}{Upper}
  }

  % reserve the space for this in a later slide.
  \invisible{%
    \ArraySwap{el87}{el45}%
  }
\end{tikzpicture}}
\vspace{3.5em}
\end{frame}

\begin{frame}[b]{Quick Sort Example}
\centerline{%
\begin{tikzpicture}
  \ArrayLabel{}
  \ArrayNode[MyGreen]{52}
  \ArrayElLabel<-12mm><0mm>[val]{pivot}
  \ArrayList[CSUgold]{13, 32, 48, 22, 11}

  \alt<1>{
    \ArrayElLabel{small}
    \ArrayNode[MyPurple]{87}
  }{
    \alt<-2>{
      \ArrayNode[CSUgold]{87}
    }{
      \ArrayNode[CSUgold]{45}
    }
    \ArrayElLabel[val]{small}
  }

  \ArrayList[MyPurple]{55, 78, 96, 58, 66, 88}



  \alt<-2>{
    \ArrayNode[]{45}
  }{
    \alt<-4>{
      \ArrayNode[]{87}
    } {
      \ArrayNode[MyPurple]{87}
    }
  }

  \ArrayElLabel{index}


  \onslide<1-2> {
    \ArraySwap{el87}{el45}
  }
  \onslide<3> {
    \ArraySwap{el45}{el87}
  }

  %\onslide<1> {

  %}

  \alt<1> {
    \ArrayGroup{el13}{el11}{Lower}
    \ArrayGroup{el87}{el88}{Upper}
  }{
    \alt<-2>{
      \ArrayGroup{el13}{el87}{Lower}
    }{
      \ArrayGroup{el13}{el45}{Lower}
    }
    \alt<-4>{
      \ArrayGroup{el55}{el88}{Upper}
    }{
      \ArrayGroup{el55}{el87}{Upper}
    }
  }
  % reserve the space for this in a later slide.
  \invisible{%
    \ArraySwap{el87}{el45}%
  }
\end{tikzpicture}}
\vspace{3.5em}
\end{frame}

\begin{frame}[b]{Quick Sort Example}
\centerline{%
\begin{tikzpicture}
  \ArrayLabel{}
  \ArrayNode[MyGreen]{52}
  \ArrayElLabel<-12mm><0mm>[val]{pivot}
  \ArrayList[CSUgold]{13, 32, 48, 22, 11, 45}
  \ArrayElLabel[val]{small}

  \ArrayList[MyPurple]{55, 78, 96, 58, 66, 88, 87}

  \onslide<2> {
    \ArraySwap{el52}{el45}
  }

  \ArrayGroup{el13}{el45}{Lower}
  \ArrayGroup{el55}{el87}{Upper}

  % reserve the space for this in a previous slide.
  \invisible{%
    \ArraySwap{el45}{el87}%
  }
\end{tikzpicture}}
\vspace{3.5em}
\end{frame}


\begin{frame}[b]{Quick Sort Example}
\centerline{%
\begin{tikzpicture}
  \ArrayLabel{}
  \ArrayList[CSUgold]{45, 13, 32, 48, 22, 11}

  \ArrayNode[MyGreen]{52}
  \ArrayElLabel[val]{pivot}

  \ArrayList[MyPurple]{55, 78, 96, 58, 66, 88, 87}

  \onslide<1> {
    \ArraySwap{el45}{el52}
  }

  \ArrayGroup{el45}{el11}{Lower}
  \ArrayGroup{el55}{el87}{Upper}

  % reserve the space for this in a previous slide.
  \invisible{%
    \ArraySwap{el52}{el87}%
  }
\end{tikzpicture}}
\pause\vspace{3.5em}
\end{frame}

% \section{Quick Sort: Implementation}

% \begin{frame}[fragile]{Quick Sort Recursive Implementation}
% \begin{lstlisting}[basicstyle=\small\ttfamily]
% // Array-based implementation
% template <typename Type>
% void quickSort(Type arr[], int size) {

%   if (size > 1) {
%     const int PIVOT_INDEX {partition(arr, size)};
%     const int UPPER_FIRST {INDEX + 1};

%     quickSort(arr, PIVOT_INDEX); // lower half
%     quickSort(arr + UPPER_FIRST, size - UPPER_FIRST);
%   }
% }
% \end{lstlisting}
% \end{frame}

% \begin{frame}[fragile]{Quick Sort Recursive Implementation}
% \begin{lstlisting}[basicstyle=\small\ttfamily]
% template <typename Type>
% int partition(Type arr[], int size) {
%   int smallIndex = 0;

%   for (int index {1}; index < size; ++index) {
%     if (arr[index] < arr[0]) { // less than pivot
%       ++smallIndex;
%       swap(arr[index], arr[smallIndex]);
%     }
%   }

%   // put the pivot in its place
%   swap(arr[0], arr[smallIndex]);
%   return smallIndex;
% }
% \end{lstlisting}
% \end{frame}

\section{Quick Sort: Analysis}

\begin{frame}{Analysis of Quick Sort}
\begin{tabular}{r l l}
\toprule
\textbf{Algorithm}        & \textbf{Comparisons} & \textbf{Swaps} \\
\midrule
  Bubble Sort    & $O(n^2)$ & $O(n^2)$ \\
  Selection Sort & $O(n^2)$ & $O(n)$ \\
  Array-Based Insertion Sort & $O(n^2)$ & $O(n^2)$ \\
  Linked-List Insertion Sort & $O(n^2)$ & $O(n)$ \\
  Quick Sort: Average Case& \onslide<2->{$O(n \log_2{n})$} &  \onslide<2->{$O(n \log_2{n})$} \\
  Quick Sort: Worst Case &  \onslide<3->{$O(n^2)$} &  \onslide<3->{$O(n^2)$} \\
\bottomrule
\end{tabular}

\begin{tabular}{r c c}
  \toprule
  \textbf{Quick Sort}        & \textbf{Comparisons } & \textbf{Swaps} \\
  \midrule
    Average Case & \onslide<2->{$(1.39) n \log_2{n}$} & \onslide<2->{$(0.69) n \log_2{n}$} \\
    Worst Case   &  \onslide<3->{$\frac{n^2 - n}{2}$} & \onslide<3->{$\frac{n^2 - 3n}{2}$} \\
  \bottomrule
\end{tabular}

% \begin{itemize}
% \item<2-> Average Case:
%   \begin{itemize}
%   \item
%     Comparisons: 
%   \item
%     Swaps: 
%   \end{itemize}
% \item<3-> Worst Case:
%   \begin{itemize}
%   \item
%     Comparisons:
%   \item
%     Swaps: 
%   \end{itemize}
% \end{itemize}

\end{frame}


\begin{frame}{Worst Case of Quick Sort}
\large\center{}
  If the chosen pivot is always the smallest or largest value, Quick Sort behaves like a \emph{Selection Sort}.

\end{frame}

\begin{frame}{Quick Sort Optimizations}
\begin{itemize}
  \item Choose a better pivot.
  \begin{itemize}
    \item Median of three: first, middle, last elements.
    \item Randomly select a pivot.
  \end{itemize}
  \item Use a different algorithm for small sub-lists.
  \begin{itemize}
    \item
      Insertion sort is best for lists of 10 or fewer elements.
    \item
      The overhead of the recursive calls is greater than the cost of the insertion sort.
  \end{itemize}
  \item Use fancy partitioning techniques.
  \begin{itemize}
    \item
      Hoare partitioning: Traverse the list from both sizes, swapping smaller elements to the left and larger elements to the right.
    \item
      Dual-Pivot partitioning: Partitions the list into three sub-lists.
  \end{itemize}
\end{itemize}
\end{frame}

\section{Merge Sort}
\begin{frame}{Merge Sort}
\begin{itemize}
\item
  \emph{Quick Sort}: $O(n\log_2n)$ average case; $O(n^2)$ worst case
\item
  \emph{Merge Sort}: always $O(n\log_2n)$

\item
  Unlike Quick Sort, it divides the list into two sub-lists of nearly equal size.
\item
  Uses the divide-and-conquer technique.
  \begin{itemize}
  \item
    Partitions the list into two sub-lists of nearly equal size.
  \item
    Sorts the sub-lists.
  \item
    Combines the sub-lists into one sorted list.
  \end{itemize}

\end{itemize}
\end{frame}


\begin{frame}{Merge Sort}
\begin{center}
  \begin{forest}
  for tree = {
    rectangle,
    color=CSUblue,
    fill=white,
    % fit=tight,
    edge=->,
    edge=very thick,
    minimum width=8mm,
    s sep = 10mm,
    l=9pt,
    font=\footnotesize
  }
  [{35\quad{}28\quad{}17\quad{}45\quad{}62\quad{}48\quad{}30\quad{}38}
    [{35\quad{}28\quad{}17\quad{}45}
      [{35\quad{}28}
        [35]
        [28]
      ]
      [{17\quad{}45}
        [17]
        [45]
      ]
    ]
    [{62\quad{}48\quad{}30\quad{}38}
      [{62\quad{}48}
        [62]
        [48]
      ]
      [{30\quad{}38}
        [30]
        [38]
      ]
    ]
  ]
  \end{forest}
\pause
  \begin{forest}
  for tree = {
    rectangle,
    grow=north,
    color=white,
    fill=CSUgold,
    % fit=tight,
    edge=<-,
    edge=very thick,
    minimum width=8mm,
    s sep = 10mm,
    l=9pt,
    for leaves={inner sep=1pt},
    font=\footnotesize
  }
  [{17\quad{}28\quad{}30\quad{}35\quad{}38\quad{}45\quad{}48\quad{}62}
    [{30\quad{}38\quad{}48\quad{}62}
      [{30\quad{}38}
        []
        []
      ]
      [{48\quad{}62}
        []
        []
      ]
    ]
    [{17\quad{}28\quad{}35\quad{}45}
      [{17 45}
        []
        []
      ]
      [{28\quad{}35}
        []
        []
      ]
    ]
  ]
  \end{forest}
\end{center}
\end{frame}

\begin{frame}{Merge Sort Algorithm}
  \large
\begin{itemize}
\item
  Uses recursion.
\item
  If the list is of  assize greater than 1,
  \begin{enumerate}
  \item
    Divide the list into two sub-lists.
  \item
    Merge sort the first sub-list.
  \item
    Merge sort the second sub-list.
  \item
    Merge the first sub-list and the second sub-list.
  \end{enumerate}
\end{itemize}
\end{frame}

\begin{frame}{Merging Two Sorted Lists}
  \begin{tikzpicture}
    \ArrayLabel{}
    \ArrayList{18, 25, 36}
    \invisible{\ArrayNode{I1}}

    \ArrayLabel[0][5]{}
    \ArrayList{20, 30, 44, 46}
    \invisible{\ArrayNode{I2}}

    \alt<-2>{
      \ArrayElLabel[el18]{left}
    } {
      \alt<-8>{
        \ArrayElLabel[el25]{left}
      }{
        \alt<-14>{
          \ArrayElLabel[el36]{left}
        }{
          \ArrayElLabel[elI1]{left}
        }
        
      }
    }
    
    \alt<-5>{
      \ArrayElLabel[el20]{right}
    }{
      \alt<-11>{
        \ArrayElLabel[el30]{right}
      }{
        \alt<-17>{
          \ArrayElLabel[el44]{right}
        }{
          \alt<-20>{
            \ArrayElLabel[el46]{right}
          }{
            \ArrayElLabel[elI2]{right}
          }
        }
      }
    }

    % Temporary Array
    \ArrayLabel[-3]{}

    % [0]
    \alt<1>{
      \ArrayNode{?}
    }{
      \ArrayNode{18}
    }
    \only<-3> {
      \ArrayElLabel{temp}
    }
    
    % [1]
    \alt<-4>{
      \ArrayNode{?}
    }{
      \ArrayNode{20}
    }
    \only<4-6> {
      \ArrayElLabel{temp}
    }

    % [2]
    \alt<-7>{
      \ArrayNode{?}
    }{
      \ArrayNode{25}
    }
    \only<7-9> {
      \ArrayElLabel{temp}
    }

    % [3]
    \alt<-10>{
      \ArrayNode{?}
    }{
      \ArrayNode{30}
    }
    \only<10-12>{
      \ArrayElLabel{temp}
    }

    % [4]
    \alt<-13>{
      \ArrayNode{?}
    }{
      \ArrayNode{36}
    }
    \only<13-15>{
      \ArrayElLabel{temp}
    }

    % [5]
    \alt<-16>{
      \ArrayNode{?}
    }{
      \ArrayNode{44}
    }
    \only<16-18>{
      \ArrayElLabel{temp}
    }

    % [6]
    \alt<-19>{
      \ArrayNode{?}
    }{
      \ArrayNode{46}
    }
    
    \only<19-21>{
      \ArrayElLabel{temp}
    }

    \invisible{\ArrayNode{}}
    \only<22->{
      \ArrayElLabel{temp}
    }
    
    \onslide<22>{}
  \end{tikzpicture}
\end{frame}

\begin{frame}{Merge Sort: Find Middle and Divide}
  \vspace{1em}
\begin{tikzpicture}[label of pointer/.append style={rotate=20}, node of list/.append style={inner sep=3pt}]
  % \LinkedList{65};
  \LinkPtr[][above]{pHead}
  \LinkedNode{65}
  \onslide<2>{\TailPtr{pMid}}
  \LinkedNode{18}
  \onslide<2>{\TailPtr{pCurr}}
  \onslide<3>{\TailPtr{pMid}}
  \LinkedNode{85}
  \onslide<4>{\TailPtr{pMid}}
  \LinkedNode{95}
  \onslide<5>{\TailPtr{pMid}}
  \onslide<3>{\TailPtr{pCurr}}
  \LinkedNode{25}
  \onslide<6->{\TailPtr{pMid}}
  \LinkedNode{20}
  \onslide<4>{\TailPtr{pCurr}}
  \LinkedNode{45}
  \LinkedNode{75}
  \onslide<5>{\TailPtr{pCurr}}
  \LinkedNode{30}
  \onslide<6>{\TailPtr[null]{pCurr}}
\end{tikzpicture}
\onslide<7->{%
  \begin{tikzpicture}[label of pointer/.append style={rotate=20}, link pointer of list/.append style={minimum height=6mm,minimum width=3mm}, node of list/.append style={inner sep=3pt}]
    \LinkPtr[][above]{pHead}
    \LinkedNode{65}
    \LinkedNode{18}
    \LinkedNode{85}
    \LinkedNode{95}
    \LinkedNode{25}
    \TailPtr{pMid}

    \LinkPtr[-1.25][above]{pFirst2}
    \LinkedNode{20}
    \LinkedNode{45}
    \LinkedNode{75}
    \LinkedNode{30}
  \end{tikzpicture}
}
\end{frame}

\begin{frame}{Merge Sort: Merge}
Sorted sub-lists are merged into one sorted list.
\begin{itemize}
\item
  Compare elements of sub-lists
\item
  Adjust pointers of nodes with the smaller datum.
\end{itemize}

\begin{tikzpicture}[
    node of list/.append style = {node distance=12.5mm, minimum width=5mm, minimum height = 7.2mm}, 
    %node of list/.append style = {inner sep=4pt},
    label of pointer/.append style={rotate=22},
    link pointer of list/.append style= { minimum width=5mm, minimum height = 7.2mm }
  ]
  %\onslide<1>{\LinkPtr{pFirst1}}
  \onslide<2->{\LinkPtr{pHead}}

  \alt<1>{
    \LinkedNode[ptrFlip1][false]{18}
  }{
    \LinkedNode[ptrFlip1]{18}
  }
  
  \only<-3>{
    \TailPtr<-4mm><4mm>{pFirst1}
  }
  \only<3-5>{\TailPtr<-4mm><-10mm>{pLastMerge}}
  
  \alt<-4>{\LinkedNode{25}}{\LinkedNode[ptrFlip3][false][elFlip2]{25}}
  \onslide<4-9>{\TailPtr<-4mm><4mm>[last]{pFirst1}}
  \onslide<9-11>{\TailPtr<-4mm><-10mm>[last]{pLastMerge}}
  \alt<-10>{\LinkedNode{65}}{\LinkedNode[ptrFlip5][false][elFlip4]{65}}
  \onslide<10-17>{\TailPtr<-4mm><8mm>[last]{pFirst1}}
  \only<17-19>{\TailPtr<-4mm><-10mm>[last]{pLastMerge}}
  \alt<-18>{\LinkedNode{85}}{\LinkedNode[ptr][false][elFlip5]{85}}
  \onslide<18->{\TailPtr<-4mm><8mm>[last]{pFirst1}}
  \LinkedNode{95}

  
  \coordinate (ptr) at (0,-2.75);
  \LinkedNode[ptr][false]{20}
  \only<-6>{\TailPtr<-4mm><4mm>{pFirst2}}
  \onslide<5->{\draw[link] (ptrFlip1) to[bend left=10] (ele);}
  \only<8->{\draw[link] (ptr) to[bend right=20] (elFlip2);}
  \onslide<6-8>{\TailPtr<-4mm><-10mm>{pLastMerge}}
  \alt<-7>{\LinkedNode{30}}{\LinkedNode[ptr][false]{30}}
  \onslide<7-12>{\TailPtr<-4mm><4mm>{pFirst2}}
  \onslide<12-13>{\TailPtr<-4mm><-10mm>{pLastMerge}}
  \only<11->{\draw[link] (ptrFlip3) to[bend left=10] (ele);}
  \LinkedNode{45}
  \onslide<13-14>{\TailPtr<-4mm><4mm>{pFirst2}}
  \onslide<14-16>{\TailPtr<-4mm><-10mm>{pLastMerge}}
  \only<16->{\draw[link] (ptr) to[bend right=10] (elFlip4);}
  \alt<-15>{\LinkedNode{75}}{\LinkedNode[ptr][false]{75}}
  \only<19->{\draw[link] (ptrFlip5) to[bend left=10] (el75);}
  \only<22->{\draw[link] (ptr) to[bend right=10] (el85);}
  \only<15->{
    \alt<-20>{
      \TailPtr<-4mm><8mm>{pFirst2}
    }{
      \TailPtr<-4mm><8mm>[null]{pFirst2}
    }
  }
  %\only<21->{\TailPtr<-4mm><4mm>[null]{pFirst2}}
  \only<20->{\TailPtr<-4mm><-10mm>{pLastMerge}}
  \onslide<22>{}
\end{tikzpicture}
\end{frame}

\section{Merge Sort: Analysis}
\begin{frame}{Analysis: Merge Sort}
\begin{itemize}
\item
  Suppose that $L$ is a list of $n$ elements, with $n > 0$.
\item
  Suppose that $n$ is a power of $2$; that is, $n = 2^m$ for some integer $m > 0$, so that we can divide the list into two sub-lists, each of size:\\
  \begin{equation*}
    \frac{n}{2}=\frac{2^m}{2}=2^{m-1}
  \end{equation*}
  $m$ will be the number of recursion levels.
\end{itemize}
\end{frame}

\begin{frame}{Analysis: Merge Sort}
\centering
\begin{forest}
  % before drawing tree={
  %     tikz+={\coordinate (a) at (current bounding box.west);},
  %   },
  % for tree={tier/.option=level},
  % level label/.style={
  %   before typesetting nodes={
  %     for nodewalk={current,tempcounta/.option=level,group={root,tree breadth-first},ancestors}{if={>OR={level}{tempcounta}}{before drawing tree={label me=##1}}{}},
  %   }
  % },
  % label me/.style={tikz+={\node [anchor=base] at (.parent |- a) {#1};}},
  block/.style={, minimum height=1.3cm},
  for tree={
    anchor=center,
    grow=south,
  },
  label/.style={anchor=center,align=left, no edge,font=\footnotesize}
  [,phantom
    [8,tier=l1
      [4,tier=l2
        [2,tier=l3
          [1,tier=l4][1,tier=l4]
        ]
        [2,tier=l3
          [1,tier=l4][1,tier=l4]
        ]
      ]
      [4,tier=l2
        [2,tier=l3
          [1,tier=l4][1,tier=l4]
        ]
        [2,tier=l3
          [1,tier=l4][1,tier=l4]
        ]
      ]
    ]
    [{Recursion Level: 0\\Calls to Merge Sort: 1 with 8 elements each}, label,tier=l1
      [{Recursion Level: 1\\Calls to Merge Sort: 2 with 4 elements each}, label,tier=l2
        [{Recursion Level: 2\\Calls to Merge Sort: 4 with 2 elements each}, label,tier=l3
          [{Recursion Level: 3\\Calls to Merge Sort: 8 with 1 elements each}, label,tier=l4]
        ]
      ]
    ]
  ]
\end{forest}
\end{frame}

\begin{frame}{Analysis: Merge Sort}
\begin{itemize}[<+->]
\item
  To merge two sorted lists of size $s$ and $t$, the maximum number of comparisons is $s + t - 1$.
\item
  Merging two sorted lists into a sorted list is where the actual comparisons and assignments are done.
\item
  Max. \# of comparisons at level $k$ of recursion:
  \begin{equation*}
    2^k\left(\frac{n}{2^k} - 1\right) = n - 2^k = O(n)
  \end{equation*}
\end{itemize}
\end{frame}

\begin{frame}{Analysis: Merge Sort}
\begin{itemize}[<+->]
\item
  There are at most \(O(n)\) comparisons per recursion level.
\item
  Therefore, there are \(O(nm)\) total comparisons, \\where \(m\) is the number of levels of recursion.
\item
  Thus, \(O(nm) \equiv O(n \log_2n)\).

\item
  Key comparisons in the worst case:\\[-1em]
    \begin{equation*}
      O\left(n\log_2n\right)
    \end{equation*}
\item
  Key comparisons in average case:\\[-1em]
    \begin{equation*}
      n\log_2n - 1.25n = O\left(n\log_2n\right)
    \end{equation*}
\end{itemize}
\end{frame}

\begin{frame}{}
\Large\centering{}
\href{https://www.toptal.com/developers/sorting-algorithms}{Sorting Algorithms Animation}
\end{frame}

\setbeamercovered{transparent}
\section{Summary}
\begin{frame}{Summary}
\begin{itemize}[<+->]
\item
  By \emph{asymptotic}, we mean the rate of growth of function $f(n)$ as $n$ grows towards infinity.
\item
  On average, a sequential search searches half the list and makes $O(n)$ comparisons, which is inefficient.
\item
  The binary-search algorithm is the optimal worst-case algorithm for solving search problems by using the \textit{comparison method}.
\item
  A binary search requires the list to be sorted but takes only $2\log_2n - 3$ or O$(\log_2n)$ key comparisons.
\end{itemize}
\end{frame}


\begin{frame}{Summary}
\begin{itemize}[<+->]
\item
  Bubble Sort: $O(n^2)$ key comparisons and item assignments.
\item
  Selection Sort: $O(n^2)$ key comparisons and $O(n)$ item assignments.
\item
  Insertion Sort: $O(n^2)$ key comparisons and item assignments.
\item
  Both the Quick Sort and Merge Sort algorithms partition a list to sort it.
  \begin{itemize}
    \item
      Quick Sort: average number of key comparisons is $O(n\log_2n)$; worst case number of key comparisons is $O(n^2)$
    \item
      Merge Sort: the number of key comparisons is $O(n\log_2n)$ on average and in the worst case.
  \end{itemize}
\end{itemize}
\end{frame}

\section*{Lab 11}
\begin{frame}{Lab 11: Insertion, Quick, and Merge Sort}
\centering\Large
  Let's take a look at Lab 11.
\end{frame}
\end{document}
